% 木星（综述）
% license CCBYSA3
% type Wiki

本文根据 CC-BY-SA 协议转载翻译自维基百科\href{https://en.wikipedia.org/wiki/Jupiter}{相关文章}。


木星是距离太阳第五颗行星，也是太阳系中体积最大的行星。它是一颗气态巨行星，其质量几乎是太阳系中所有其他行星总和的 2.5 倍，但仍不到太阳质量的千分之一。它的直径是地球的 11 倍，是太阳的十分之一。木星以 5.20 天文单位（778.5 吉米）的平均距离绕太阳运行，公转周期为 11.86 年。在地球夜空中，木星是仅次于月球和金星的第三亮的天然天体，自史前时代以来便为人类所观察。它的名字源自古罗马宗教中至高神朱庇特（Jupiter）。

木星是太阳行星中最先形成的一颗，其在太阳系原始阶段向内迁移的过程深刻影响了其他行星的形成历史。木星的大气按质量计由 76\% 的氢和 24\% 的氦组成，内部密度更高。它还含有微量的碳、氧、硫、氖等元素，以及氨、水蒸气、磷化氢、硫化氢和烃类等化合物。木星的氦含量约为太阳的 80\%，与土星的成分相似。

木星外层大气被分为一系列纬向云带，这些云带在交界处产生湍流和风暴；其中最显著的结果就是“大红斑”——自 1831 年以来被记录的一场巨型风暴。由于木星自转极快（约 10 小时自转一周），使其呈现为一颗扁球体；与两极相比，赤道部位存在约 6.5\% \(^\text{[e]}\)的明显隆起。科学家认为，木星内部结构包括一层由流体金属氢组成的外地幔，以及由更高密度物质构成的弥散内核。木星内部持续的收缩过程释放出比其从太阳接收的还要多的热量。木星的磁场是太阳系中最强且连续范围第二大的结构，由流体金属氢内部的涡电流产生。太阳风与其磁层相互作用，使磁层向外延伸并影响木星的运行环境。

至少有 97 颗卫星绕木星运行；其中四颗最大的卫星——木卫一（Io）、木卫二（Europa）、木卫三（Ganymede）和木卫四（Callisto）——位于其磁层内，用普通双筒望远镜即可观测。四者中最大的木卫三比行星水星还要大。木星周围还存在一个暗淡的行星环系统。木星环主要由尘埃组成，分为三部分：内侧呈光环状尘粒环的“晕环”（halo）、相对明亮的主环（main ring），以及外侧的薄雾状“薄环”（gossamer ring）。这些行星环在可见光和近红外光下呈现红色。木星环系统的年龄尚不清楚，可能追溯至木星形成时期。自 1973 年以来，已有九个机器人探测器访问过木星：包括七次飞掠任务和两艘专门的轨道探测器（另有两艘在途中）。在其他行星系统中也已发现类似木星的系外行星。
\subsection{名称与符号}.
Jupiter 在古希腊和古罗马文明中都以神系中的最高神命名：对希腊人而言是宙斯（Zeus），对罗马人而言是朱庇特（Jupiter）\(^\text{[20]}\)。国际天文学联合会于 1976 年正式采用 Jupiter 作为该行星的名称，并自此将新发现的卫星以这位神祇的情人、宠爱者及后代命名\(^\text{[21]}\)。木星的行星符号 ♃ 源自带有水平笔画的希腊字母 zeta（⟨Ƶ⟩），作为 Zeus 的缩写\(^\text{[22][23]}\)。在拉丁语中，Iovis 是 Iuppiter（即 Jupiter）的属格形式，并与 Zeus（“天空之父”）的词源相关。英语中的对应形式 Jove 被认为在 14 世纪起作为该行星的诗性名称出现\(^\text{[24]}\)。Jovian 是 Jupiter 的形容词形式。较早的形容词 jovial，被中古时代的占星家使用，后来演变为“愉快的”或“欢乐的”含义，此情绪被认为与木星在占星学中的影响相关\(^\text{[25]}\)。原始希腊神祇 Zeus 还提供了词根 zeno-，用于构成一些与木星相关的词语，例如 zenography\(^\text{[f]}\)。
\subsection{形成与迁移}
木星被认为是太阳系中最古老的行星，其形成时间仅在太阳之后一百万年，约早于地球五千万年\(^\text{[26]}\)。当前的太阳系形成模型表明，木星形成于雪线处或其以外的位置：即距离早期太阳足够远、温度足够低，使水等挥发物能够凝结成固体的区域\(^\text{[27]}\)。木星首先形成固体核心，随后吸积其气态大气层。因此，该行星必须在太阳星云完全消散之前形成\(^\text{[28]}\)。在形成过程中，木星的质量逐渐增加，直到达到地球质量的 20 倍，其中约一半由硅酸盐、冰以及其他重元素成分构成\(^\text{[26]}\)。当原始木星增大超过 50 个地球质量时，它在太阳星云中开辟出一个空隙\(^\text{[26]}\)。此后，这颗不断成长的行星在 300–400 万年内达到最终质量\(^\text{[26][28]}\)。

根据“Grand tack 假说”，木星最初在距太阳约 3.5 AU（5.20 亿公里；3.30 亿英里）的位置开始形成。随着年轻的木星吸积质量，它与环绕太阳的气体盘相互作用，并受到土星轨道共振的影响，从而向内迁移\(^\text{[27][29]}\)。这扰乱了数颗在更靠近太阳轨道上运行的超级地球，使它们发生破坏性的碰撞\(^\text{[30]}\)。随后，土星开始以比木星更快的速度向内迁移，直到两者在距太阳约 1.5 AU（2.20 亿公里；1.40 亿英里）处被俘获进入 3:2 平均运动共振\(^\text{[31]}\)。这一事件改变了迁移方向，使它们开始远离太阳并迁出内太阳系，移动至如今的位置\(^\text{[30]}\)。上述全过程持续约 300–600 万年，而木星最终阶段的迁移历经数十万年完成\(^\text{[29][32]}\)。木星从内太阳系的迁移最终使包括地球在内的内行星得以从残余碎片中形成\(^\text{[33]}\)。

Grand tack 假说仍存在多个未解决的问题。其推导出的类地行星形成时间尺度似乎与测量得到的元素组成不一致\(^\text{[34]}\)。如果木星曾穿越太阳星云迁移，它最终可能会定居在更靠近太阳的轨道上\(^\text{[35]}\)。一些太阳系形成的竞争模型预测木星的形成轨道性质与今日行星的位置非常接近\(^\text{[28]}\)。其他模型则预测木星形成于更为外侧的位置，例如 18 AU（27 亿公里；17 亿英里）\(^\text{[36][37]}\)。

根据 Nice 模型，在太阳系历史最初的 6 亿年间，原始柯伊伯带天体的内落导致木星和土星从它们最初的位置迁移至 1:2 共振，这使土星移入更高轨道，从而扰乱了天王星和海王星的轨道，消耗了柯伊伯带，并触发了晚期重轰炸\(^\text{[38]}\)。

根据 Jumping-Jupiter 情景，木星在早期太阳系中的迁移可能导致第五颗气态巨行星被抛出。该假说认为，在其轨道迁移过程中，木星的引力扰动破坏了其他气态巨行星的轨道，可能使其中一颗完全被逐出太阳系。此类事件的动力学将极大改变太阳系的形成与配置，最终只留下如今人类所见的四颗气态巨行星\(^\text{[39]}\)。

基于木星的组成，有研究者提出其最初形成于分子氮（N\(_2\)）雪线之外，即距太阳 20–30 AU（30–45 亿公里；19–28 亿英里）的位置，甚至可能形成于氩雪线之外，后者可能远至 40 AU（60 亿公里；37 亿英里）\(^\text{[40][41]}\)。如果木星形成于这些极远的位置之一，则它将在大约 70 万年的时间尺度上向内迁移至目前的位置\(^\text{[36][37]}\)，这一过程发生于行星开始形成后约 200–300 万年的时代。在该模型中，土星、天王星和海王星的形成位置将比木星更远，而土星亦会向内迁移\(^\text{[36]}\)。
\subsection{物理特征}
木星是一颗气态巨行星，这意味着其化学组成主要为氢和氦。这些材料在行星地质学中被分类为“气体”，这一术语并不指示其物质状态。木星是太阳系中最大的行星，其赤道直径为 142,984 公里（88,846 英里），体积为地球的 1,321 倍\(^\text{[3][42]}\)。其平均密度为 1.326 g/cm\(^3\)\(^\text{[g]}\)，低于四颗类地行星的密度\(^\text{[44][45]}\)。
\subsubsection{成分}
木星大气按质量计约由 76\% 的氢和 24\% 的氦组成。按体积计，上层大气约含 90\% 的氢和 10\% 的氦，这一较低比例是因为单个氦原子比该大气层中形成的氢分子更为致密\(^\text{[46]}\)。大气含有微量的碳、氧、硫和氖\(^\text{[47]}\)，以及氨、水蒸气、磷化氢、硫化氢和甲烷、乙烷、苯等烃类化合物\(^\text{[48]}\)。其最外层包含冻结氨的晶体\(^\text{[49]}\)。行星内部密度更高，按质量计约由 71\% 的氢、24\% 的氦和 5\% 的其他元素组成\(^\text{[50][51]}\)。

大气中氢与氦的比例接近理论上的原始太阳星云成分\(^\text{[52]}\)。上层大气中的氖含量为每百万份质量中的 20 份，约为太阳丰度的十分之一\(^\text{[53]}\)。木星的氦含量约为太阳的 80\%，这是由于这些元素在行星内部深处以富含氦的液滴形式沉降的过程所致\(^\text{[54][55]}\)。

基于光谱学分析，土星的组成被认为与木星相似，但另外两颗巨行星天王星和海王星含有相对更少的氢和氦，而相对含有更多的次常见元素，包括氧、碳、氮和硫\(^\text{[56]}\)。这些行星被称为冰巨行星，因为在其形成期间，这些元素被认为以冰的形式并入其中；然而它们内部实际上可能含有极少量的冰\(^\text{[57]}\)。
\subsubsection{大小与质量}
\begin{figure}[ht]
\centering
\includegraphics[width=8cm]{./figures/82e5bc11e94dd9a1.png}
\caption{木星与地球及地球月球的尺寸比较} \label{fig_Jupite_1}
\end{figure}
木星的直径约为地球的十一倍（11.209 R🜨）；其质量是地球的 318 倍\(^\text{[3]}\)，相当于太阳系中其他所有行星质量总和的 2.5 倍。其质量之大，以至于它与太阳的质心位于太阳表面之外，距太阳中心 1.068 个太阳半径\(^\text{[58][59]}\)。木星的半径约为太阳半径的十分之一（0.10276 R☉）\(^\text{[60]}\)，其质量约为太阳质量的千分之一，两者的平均密度相似\(^\text{[61]}\)。“木星质量”（MJ 或 MJup）常被用作描述其他天体质量的单位，特别是系外行星和棕矮星。例如，系外行星 HD 209458 b 的质量为 0.69 MJ，而棕矮星 Gliese 229 b 的质量为 60.4 MJ\(^\text{[62][63]}\)。

木星辐射的热量多于其从太阳接收的辐射能量，这是因为其内部收缩过程中的 Kelvin–Helmholtz 机制\(^\text{[64]:30^\text{}[65]}\)。这一过程导致木星每年缩小约 1 毫米（0.039 英寸）\(^\text{[66][67]}\)。在其形成之初，木星更为炽热，且其直径约为当前的两倍\(^\text{[68][69]}\)。

理论模型表明，如果木星再增加 40\% 以上的质量，其内部将因压缩而导致体积缩小，尽管质量增加。对较小的质量变化而言，其半径变化并不显著\(^\text{[70]}\)。因此，木星被认为已达到由其组成与演化历史所能达到的最大行星直径\(^\text{[71]}\)。随着质量进一步增加而导致的收缩过程会持续，直至达到明显的恒星点火为止\(^\text{[72]}\)。尽管木星需要约 75 倍目前质量才能点燃氢聚变、成为一颗恒星\(^\text{[73]}\)，但其直径依然符合最小红矮星可能仅略大于土星半径的事实\(^\text{[74]}\)。
\subsubsection{内部结构}
\begin{figure}[ht]
\centering
\includegraphics[width=8cm]{./figures/2dfe73e2dbd3a048.png}
\caption{木星的内部结构、表面特征、行星环和内侧卫星示意图} \label{fig_Jupite_2}
\end{figure}
在 21 世纪初之前，大多数科学家提出了两种关于木星形成的主要情景之一。如果行星首先作为固体天体吸积形成，它将由一个致密核心构成，其外包围一层流体金属氢（含少量氦），这一层向外延伸至行星半径的约 80\%\(^\text{[75]}\)，并有一层主要由分子氢组成的外层大气\(^\text{[67]}\)。另一种可能是，如果行星直接从气态原行星盘中坍缩形成，则预计它将完全缺乏核心，而是从中心向外由逐渐变得更致密的流体构成（主要是分子氢与金属氢）。Juno 任务的数据表明，木星具有一个向地幔混合的弥散核心，延伸至行星半径的 30–50\%，由总质量相当于 7–25 个地球质量的重元素构成\(^\text{[76][77][78]}\)。这种混合过程可能在形成期间便已发生，当时行星从周围星云中吸积固体与气体\(^\text{[79]}\)。另一种可能是，在木星形成几百万年后，一颗约十倍地球质量的行星发生碰撞，扰乱了原本致密的木星核心\(^\text{[77][80]}\)。

在金属氢层之外，是透明的内部氢大气。在这一深度，压力和温度高于分子氢的临界压力 1.3 MPa 与临界温度 33 K（−240.2 °C；−400.3 °F）\(^\text{[81]}\)。在这种状态下，不再存在明确的液态与气态的分界——氢处于所谓的超临界流体状态。从云层向下延伸的氢气和氦气在更深处逐渐过渡为液体，可能类似于由液态氢及其他超临界流体构成的“海洋”\(^\text{[64]:22^\text{}[82][83]}\)。从物理上讲，随着深度增加，气体会逐渐变得更热、更致密\(^\text{[84][85]}\)。

氦和氖如同“雨滴”般在大气下层向下沉降，使上层大气中的这些元素含量减少\(^\text{[54][86]}\)。计算表明，氦滴在距中心 60,000 公里（37,000 英里）（即云顶下方 11,000 公里［6,800 英里］）处从金属氢中分离，并在 50,000 公里（31,000 英里）（云层下方 22,000 公里［14,000 英里］）处重新混合\(^\text{[87]}\)。研究指出，如同土星\(^\text{[88]}\)以及冰巨行星天王星和海王星\(^\text{[89]}\)一样，木星内部也可能存在钻石雨。

木星内部的温度和压力向内稳步升高，因为行星形成时的热量只能通过对流方式散失\(^\text{[55]}\)。在大气压力约为一个地球标准大气压（约 0.10 MPa（1 bar））的深度，温度约为 165 K（−108 °C；−163 °F）。超临界氢从分子流体逐渐过渡到金属流体的区域涵盖 50–400 GPa 的压力范围，对应的温度为 5,000–8,400 K（4,730–8,130 °C；8,540–14,660 °F）。木星弥散核心的温度估计为 20,000 K（19,700 °C；35,500 °F），压力约为 4,000 GPa\(^\text{[90]}\)。
\subsubsection{大气}
木星大气主要由分子氢和氦组成，含有较少量的水、甲烷、硫化氢和氨等其他化合物\(^\text{[91]}\)。木星大气自云层向下延伸约 3,000 公里（2,000 英里）\(^\text{[90]}\)。

\textbf{云层}
\begin{figure}[ht]
\centering
\includegraphics[width=8cm]{./figures/a270704535c387fb.png}
\caption{木星云层系统在一个月中的运动延时影像（1979 年旅行者 1 号飞掠期间拍摄）} \label{fig_Jupite_3}
\end{figure}
木星上始终覆盖着由氨晶体构成的云层，其中可能还含有硫氢化铵\(^\text{[92]}\)。这些云位于大气的对流层顶，形成分布在不同行星纬度上的云带，被称为热带区域。它们进一步被划分为较浅色的“区”（zones）与较深色的“带”（belts）。这些相互冲突的环流模式的相互作用引发风暴与湍流。在纬向急流中，常见的风速可达每秒 100 米（360 公里/小时；220 英里/小时）\(^\text{[93]}\)。这些云带的宽度、颜色与强度会因年份而变化，但其整体结构足够稳定，使科学家能够为它们命名\(^\text{[59]:6^\text{}]}\)。

云层厚约 50 公里（31 英里），由至少两层氨云构成：顶部为一层较薄、较透明的区域，下方为更厚的下层云带。可能在氨云之下还存在一薄层水云，这一推测来自对木星大气中闪电现象的探测\(^\text{[94]}\)。这些放电的能量可高达地球闪电的千倍\(^\text{[95]}\)。水云被认为以与地球雷暴相似的方式形成雷暴，由木星内部上升的热量驱动\(^\text{[96]}\)。Juno 任务揭示了“浅层闪电”的存在，其起源于大气中较高高度的氨-水云\(^\text{[97]}\)。这些放电会携带由水-氨泥浆构成、外覆冰壳的“浆球”（mushballs），它们会坠入大气深处\(^\text{[98]}\)。在木星高层大气中还观测到类似地球“精灵”与“红色妖精”的上层大气闪电——持续约 1.4 毫秒的亮光，由于氢元素的缘故呈现蓝色或粉红色\(^\text{[99][100]}\)。

木星云层中的橙色与棕色由下沉的化合物造成，这些化合物在受到来自太阳的紫外光照射后发生变色。其具体成分尚不确定，但普遍认为由磷、硫或可能的烃类构成\(^\text{[64]:39^\text{}[101]}\)。这些被称为“着色团”的化合物与来自下层的暖云混合。浅色云区则是因上升对流产生氨晶体，这些晶体将着色团遮蔽于视线之外\(^\text{[102]}\)。

木星自转轴倾角很小，因此其两极永远接收的太阳辐射比赤道区域少。行星内部的对流将能量输送至两极，从而在云层高度维持热量平衡\(^\text{[59]:54^\text{}]}\)。

\textbf{大红斑及其他涡旋}
\begin{figure}[ht]
\centering
\includegraphics[width=8cm]{./figures/b90e5abcdffbaa73.png}
\caption{由朱诺号探测器以真实色彩拍摄的大红斑特写。由于朱诺号的成像方式，拼接后的图像呈现出显著的桶形畸变。} \label{fig_Jupite_4}
\end{figure}
木星最著名的特征之一是大红斑\(^\text{[103]}\)，这是一场位于南纬 22° 的持久反气旋风暴。它首次于 1831 年被观测到\(^\text{[104]}\)，甚至可能早在 1665 年就已被记录\(^\text{[105][106]}\)。哈勃太空望远镜所获得的图像显示，在大红斑附近还有另外两个“红斑”\(^\text{[107][108]}\)。通过口径 12 厘米或更大的地基望远镜即可观测到这场风暴\(^\text{[109]}\)。风暴呈逆时针方向旋转，其周期约为六天\(^\text{[110]}\)。该风暴的最高高度约比周围云顶高出 8 公里（5 英里）\(^\text{[111]}\)。大红斑的具体成分及其红色来源仍不确定，不过光解氨与乙炔发生反应是可能的解释之一\(^\text{[112]}\)。

大红斑的面积大于地球\(^\text{[113]}\)。数学模型表明，这场风暴稳定存在，将成为木星上的长期特征\(^\text{[114]}\)。然而，从发现至今它的尺寸已显著缩小。19 世纪末的早期观测显示其宽度约为 41,000 公里（25,500 英里）。截至 2015 年，风暴的大小约为 16,500 × 10,940 公里（10,250 × 6,800 英里）\(^\text{[115]}\)，且其长度每年缩减约 930 公里（580 英里）\(^\text{[113]}\)。2021 年 10 月，Juno 飞掠任务测量了大红斑的深度，结果约为 300–500 公里（190–310 英里）\(^\text{[116]}\)。

Juno 任务在木星两极发现了数个气旋群。北极气旋群包含九个气旋，其中一个位于中心，周围分布着另外八个；南极气旋群同样有一个中心涡旋，但周围由五个大型风暴与一个较小的风暴组成，总计七个\(^\text{[117][118]}\)。

在 2000 年，木星南半球出现了一个类似大红斑但规模较小的大气特征。这是由几个较小的白色椭圆风暴合并而成——这三个较小白色椭圆最初形成于 1939–1940 年。合并后的结构被命名为 Oval BA。此后，它的强度有所增强，颜色也由白色转为红色，因此获得了“小红斑”之称\(^\text{[119][120]}\)。

2017 年 4 月，在木星北极的热层中发现了一个“大冷斑”。这一结构宽 24,000 公里（15,000 英里），高 12,000 公里（7,500 英里），其温度比周围材料低 200 °C（360 °F）。虽然该冷斑在短期内形态与强度会变化，但其在大气中的总体位置已保持超过 15 年。它可能是一种类似大红斑的巨大涡旋，并似乎像地球热层中的涡旋一样呈准稳定状态。该结构可能由来自木卫一的带电粒子与木星强磁场相互作用形成，导致热量流动被重新分配\(^\text{[121]}\)。
\subsubsection{磁层}
\begin{figure}[ht]
\centering
\includegraphics[width=8cm]{./figures/3a2e5013a5750a65.png}
\caption{伽利略卫星对木星磁层的影响} \label{fig_Jupite_5}
\end{figure}
木星的磁场是太阳系所有行星中最强的\(^\text{[102]}\)，其偶极矩为 4.170 高斯（0.4170 mT），并相对于自转轴倾斜 10.31°。其表面磁场强度从 2 高斯（0.20 mT）到 20 高斯（2.0 mT）不等\(^\text{[122]}\)。该磁场被认为由流体金属氢核心内部的涡电流——导电物质的旋转运动——所产生。在距离木星约 75 个木星半径处，磁层与太阳风的相互作用产生弓形激波。环绕木星磁层外侧的是磁暂停层，其位置位于磁鞘层的内缘——这是位于磁暂停层与弓形激波之间的区域。太阳风与这些区域相互作用，使磁层在木星背风侧被拉长并向外延伸，几乎抵达土星轨道。木星最大的四颗卫星全部在磁层内运行，从而受到保护免受太阳风影响\(^\text{[64]:69^\text{}]}\)。

木卫一的火山喷发大量释放二氧化硫，在其轨道附近形成一个气体环。气体在木星磁层中被电离，产生硫和氧离子。它们与来自木星大气的氢离子一起，在木星赤道平面形成一个等离子体片。该等离子体片随着行星共同旋转，使偶极磁场被拉伸为磁盘状。等离子体片中的电子产生强烈的无线电信号，其中叠加着 0.6–30 MHz 范围内的短促爆发，这些信号可通过地面消费级短波无线电接收机探测到\(^\text{[123][124]}\)。当木卫一穿过该气体环时，其相互作用会产生阿尔芬波，将电离物质带入木星极区。因此，通过回旋脉塞机制产生无线电波，能量沿锥形表面传播。当地球穿过该锥面时，来自木星的无线电辐射可超过太阳的无线电输出\(^\text{[125]}\)。
\subsubsection{行星环}
\begin{figure}[ht]
\centering
\includegraphics[width=8cm]{./figures/143fda39a5f08d29.png}
\caption{以红外光拍摄的木星图像，显示其微弱的行星环，以及两颗卫星——Amalthea 与 Adrastea、极光及大气特征。} \label{fig_Jupite_6}
\end{figure}
木星拥有一个由三部分组成的微弱行星环系统：最内侧由颗粒构成的光环状“晕环”（halo），中间较为明亮的主环，以及最外侧的薄雾环（gossamer ring）\(^\text{[126]}\)。这些行星环似乎由尘埃组成，而土星的行星环则由冰构成\(^\text{[64]:65^\text{}]}\)。主环极可能由卫星 Adrastrea 与 Metis 被撞击后喷射出的物质构成，这些物质因木星强大的引力影响而被吸引进入行星环中。新的物质则通过额外撞击持续补充\(^\text{[127]}\)。类似地，Thebe 与 Amalthea 两颗卫星被认为产生了薄雾环中的两个不同尘埃成分\(^\text{[127]}\)。有证据显示存在第四条行星环，可能由 Amalthea 碰撞产生的碎片构成，并沿着该卫星的轨道分布\(^\text{[128]}\)。
\subsection{轨道与自转}
\begin{figure}[ht]
\centering
\includegraphics[width=8cm]{./figures/9931f8d9015015fb.png}
\caption{展示木星自转及其卫星轨道运动的三小时延时影像} \label{fig_Jupite_7}
\end{figure}
木星是唯一一颗其与太阳的质心位于太阳体积之外的行星，不过仅超过太阳半径的 7\%\(^\text{[129][130]}\)。木星与太阳之间的平均距离为 7.78 亿公里（5.20 AU），其绕太阳公转一周需 11.86 年。这大约是土星轨道周期的五分之二，构成一个近似的轨道共振\(^\text{[131]}\)。木星轨道平面相对于地球倾斜 1.30°。由于其轨道偏心率为 0.049，木星在近日点时比在远日点时更接近太阳约 7,500 万公里\(^\text{[3]}\)，意味着其轨道几乎是圆形的。与此相比，系外行星的发现揭示许多木星大小的行星具有高轨道偏心率。模型表明，这可能是因为太阳系中存在两颗巨行星；而拥有三颗或更多巨行星的系统更倾向于产生较大的偏心率\(^\text{[132]}\)。

木星的自转轴倾角为 3.13°，相对较小，因此其季节变化远不及地球与火星显著\(^\text{[133]}\)。

木星的自转速度是太阳系所有行星中最快的，自转周期略短于十小时；这导致其赤道隆起，通过业余望远镜便可轻松观测。由于木星并非固体，其上层大气呈现差异自转。木星极区大气的自转速度比赤道大气慢约五分钟\(^\text{[134]}\)。木星是一颗扁球体，这意味着其赤道直径比两极之间的直径更长\(^\text{[85]}\)。在木星上，赤道直径比极径长 9,276 公里（5,764 英里）\(^\text{[3]}\)。

跟踪行星自转时，尤其是在绘制大气特征运动时，会使用三个参考系统。系统 I 适用于 7°N 至 7°S 之间的纬度，其自转周期为行星最短，为 9 小时 50 分 30.0 秒。系统 II 适用于其余纬度，自转周期为 9 小时 55 分 40.6 秒\(^\text{[135]}\)。系统 III 由射电天文学家定义，对应于木星磁层的自转；其周期被视为木星的官方自转周期\(^\text{[136]}\)。
\subsection{观测}
\begin{figure}[ht]
\centering
\includegraphics[width=8cm]{./figures/a02a6f8994823765.png}
\caption{通过业余望远镜观测到的木星及其四颗伽利略卫星} \label{fig_Jupite_8}
\end{figure}
木星通常是夜空中第四亮的天体（仅次于太阳、月亮和金星）\(^\text{[102]}\)，尽管在冲日附近火星有时会比木星更亮。根据木星相对于地球的位置不同，其视星等可在冲日时亮至 −2.94，而在与太阳合相时暗至 −1.66\(^\text{[17]}\)。其平均视星等为 −2.20，标准差为 0.33\(^\text{[17]}\)。木星的视直径也随之变化，在 50.1 至 30.5 角秒之间\(^\text{[3]}\)。当木星经过轨道近日点时会出现有利的冲日，使其更接近地球\(^\text{[137]}\)。在冲日前后，木星会出现大约 121 天的逆行运动，在天空中逆向移动约 9.9°，之后恢复顺行\(^\text{[138]}\)。

由于木星的轨道位于地球轨道之外，从地球观测到的木星相位角始终小于 11.5°；因此，通过地基望远镜看到的木星总是几乎完全被照亮的。只有在飞掠木星的航天器任务中，才获得了呈弦月状的木星影像\(^\text{[139]}\)。在小型望远镜下通常可以看到木星的四颗伽利略卫星，以及横跨木星大气的云带。口径约 4–6 英寸（10–15 厘米）的较大望远镜在大红斑朝向地球时能够观测到该结构\(^\text{[140][141]}\)。
\subsubsection{历史}
\textbf{望远镜出现前的研究}
\begin{figure}[ht]
\centering
\includegraphics[width=8cm]{./figures/b20ffb6befc24739.png}
\caption{《至大论》（Almagest）中木星（☉）相对于地球（🜨）经度运动的模型} \label{fig_Jupite_9}
\end{figure}
对木星的观测在公元前 7–8 世纪的巴比伦天文学家记录中已有确证\(^\text{[142]}\)。古代中国人称木星为“岁星”（Suìxīng 歲星），并依据木星绕太阳一周约需十二年左右的时间建立了十二地支体系；汉语至今仍以“岁”字指代年龄的“年”。到了公元前 4 世纪，这些观测发展为中国的黄道十二宫\(^\text{[143]}\)，每一年都与一位“太岁”星和掌管木星在夜空位置相对区域的天神相关联。这些信仰在部分道教与民间宗教实践以及东亚黄道的十二生肖中仍有保留。中国历史学家席泽宗提出，古代天文学家甘德\(^\text{[144]}\)曾报告过一颗与木星“相伴”的小星\(^\text{[145]}\)，这可能意味着他曾以肉眼观测到木星的一颗卫星。如果属实，这一发现将比伽利略的观测早近两千年\(^\text{[146][147]}\)。

2016 年的一篇论文指出，巴比伦人在公元前 50 年以前便使用梯形法则来对木星在黄道上的速度进行积分\(^\text{[148]}\)。在公元 2 世纪的著作《至大论》（Almagest）中，希腊化时期的天文学家托勒密（Claudius Ptolemaeus）构建了一个基于本轮与均轮的地心行星模型，用以解释木星相对于地球的运动，并给出其绕地球公转周期为 4332.38 天，即 11.86 年\(^\text{[149]}\)。

\textbf{地基望远镜研究}
\begin{figure}[ht]
\centering
\includegraphics[width=8cm]{./figures/9ce1392db0ded268.png}
\caption{伽利略在《星际信使》（Sidereus Nuncius）中绘制的木星及其“美第奇星”（Medicean Stars）图像} \label{fig_Jupite_10}
\end{figure}
1610 年，意大利博学者伽利略·伽利莱（Galileo Galilei）使用望远镜发现了木星的四颗最大卫星（现称伽利略卫星）。这被认为是首次利用望远镜观测地球以外的卫星。在伽利略之后仅一天，西蒙·马留斯（Simon Marius）也独立发现了木星的卫星，尽管他直到 1614 年才在书中发表其发现\(^\text{[150]}\)。然而沿用至今的是马留斯为四颗主要卫星所命名的名称：Io、Europa、Ganymede 和 Callisto。此项发现为尼古拉·哥白尼的日心说行星运动理论提供了强有力的支持；伽利略对日心说的公开支持导致他被宗教裁判所审判并定罪\(^\text{[151]}\)。

1639 年秋，来自那不勒斯的光学技师弗朗切斯科·丰塔纳（Francesco Fontana）测试了他自制的一台 22 棕榈长的望远镜，并发现了木星大气的特征性云带\(^\text{[152]}\)。

1660 年代，乔瓦尼·卡西尼（Giovanni Cassini）使用新型望远镜发现了木星大气中的斑点，观察到木星呈扁球形，并估算了其自转周期\(^\text{[153]}\)。1692 年，卡西尼注意到木星大气存在差异自转现象\(^\text{[154]}\)。

大红斑可能早在 1664 年由罗伯特·胡克（Robert Hooke）观测到，1665 年由卡西尼观测到，但这存在争议。药剂师海因里希·施瓦贝（Heinrich Schwabe）于 1831 年绘制了已知最早展示大红斑细节的图像\(^\text{[155]}\)。据记录，大红斑在 1665 年至 1708 年间数次从视野中消失，直到 1878 年才再次变得十分显眼\(^\text{[156]}\)。它在 1883 年以及 20 世纪初又被记录为再次变暗\(^\text{[157]}\)。

乔瓦尼·博雷利（Giovanni Borelli）和卡西尼均为木星卫星的运动制作了精确的表格，使人们能够预测卫星何时会从木星前方或后方经过。到 1670 年代，卡西尼注意到当木星位于太阳与地球的对侧时，这些现象发生的时间会比预计晚约 17 分钟。奥勒·罗默（Ole Rømer）据此推断光并非瞬时传播（这一结论先前被卡西尼否认）\(^\text{[51]}\)，并利用这一时间差估算了光速\(^\text{[158][159]}\)。

1892 年，E. E. 巴纳德（E. E. Barnard）利用加州利克天文台 36 英寸（910 毫米）折射望远镜观测到木星的第五颗卫星。此卫星后来被命名为 Amalthea\(^\text{[160]}\)。这是最后一颗通过望远镜视觉观测直接发现的行星卫星\(^\text{[161]}\)。在 1979 年旅行者 1 号飞掠之前，又有另外八颗木星卫星被发现\(^\text{[h]}\)。

1932 年，鲁珀特·维尔特（Rupert Wildt）在木星光谱中识别出氨与甲烷的吸收带\(^\text{[162]}\)。1938 年，木星大气中发现了三个长期存在的反气旋特征，被称为“白椭圆”。几十年来，这些椭圆在彼此靠近时并未发生合并。最终，其中两个椭圆于 1998 年合并，随后在 2000 年吸收了第三个，形成了 Oval BA\(^\text{[163]}\)。

\textbf{射电望远镜研究}

1955 年，伯纳德·伯克（Bernard Burke）和肯尼斯·富兰克林（Kenneth Franklin）发现木星会以 22.2 MHz 的频率发射无线电脉冲\(^\text{[64]:36^\text{}]}\)。这些脉冲的周期与木星的自转周期一致，他们利用这一信息更精确地测定了木星的自转速率。研究发现，来自木星的无线电脉冲有两种形式：持续数秒的长脉冲（L-bursts）以及持续不足百分之一秒的短脉冲（S-bursts）\(^\text{[164]}\)。

科学家已发现来自木星的三类无线电信号：
\begin{itemize}
\item 十米波无线电脉冲（波长达数十米）随木星自转变化，并受到木卫一与木星磁场相互作用的影响\(^\text{[165]}\)。
\item 分米波无线电辐射（波长以厘米计）首次由弗兰克·德雷克（Frank Drake）和海因·赫瓦图姆（Hein Hvatum）于 1959 年观测到\(^\text{[64]:36^\text{}]}\)。这种信号来自木星赤道周围的一个环状带，带中的电子在木星磁场的加速下产生回旋辐射\(^\text{[166]}\)。
\item 热辐射由木星大气中的热量产生\(^\text{[64]:43^\text{}]}\)。
\end{itemize}
\textbf{探测}
自 1973 年以来，自动化航天器开始探访木星，当时先驱者 10 号（Pioneer 10）探测器首次飞掠木星并传回了关于其性质和现象的突破性数据\(^\text{[167][168]}\)。前往木星的任务需要付出能量代价，其大小通常由航天器所需的速度增量（delta-v）衡量。从近地轨道进入通往木星的霍曼转移轨道需要 6.3 km/s 的速度增量\(^\text{[169]}\)，这一数值与到达近地轨道所需的 9.7 km/s 相当\(^\text{[170]}\)。通过行星飞掠提供的引力助推可以降低抵达木星所需的能量\(^\text{[171]}\)。

\textbf{飞掠任务}
\begin{figure}[ht]
\centering
\includegraphics[width=8cm]{./figures/ef4b56f8efb8c77d.png}
\caption{} \label{fig_Jupite_11}
\end{figure}
自 1973 年起，多艘航天器进行了行星飞掠机动，使其进入木星的可观测范围。先驱者号任务首次获得了木星大气及其多颗卫星的近距离图像。它们发现木星附近的辐射场远比预期强，但两艘探测器都成功在该恶劣环境中生存下来。这些航天器的轨道被用来进一步精确木星系统的质量估计。通过行星造成的无线电掩星现象，使得对木星直径及两极扁率的测量更加精确\(^\text{[59]:47^\text{}[173]}\)。

六年后，旅行者号任务极大提升了对伽利略卫星的认识，并发现了木星的行星环。它们也确认大红斑是一场反气旋风暴。图像对比显示，自先驱者号任务以来，大红斑的色调已经发生变化，从橙色转为深棕色。沿着木卫一轨道路径发现了一个由电离原子构成的环，其来源被确认是木卫一表面爆发的火山活动。当航天器从木星背面飞过时，它观测到夜侧大气中的闪电\(^\text{[59]:87^\text{}[174]}\)。

下一次遭遇木星的任务是尤利西斯号（Ulysses）太阳探测器。1992 年 2 月，它进行了一次飞掠机动以进入太阳的极轨道。在此次飞掠中，探测器研究了木星的磁层，但由于未配备摄像设备，因此无法拍摄行星图像。六年后，该探测器再次经过木星附近，但距离远得多\(^\text{[172]}\)。

2000 年，卡西尼号（Cassini）探测器在前往土星途中飞掠木星，并提供了分辨率更高的图像\(^\text{[175]}\)。

2007 年，新视野号（New Horizons）探测器在前往冥王星途中飞掠木星以获得引力助推\(^\text{[176]}\)。探测器的摄像设备测量了来自木卫一火山的等离子体输出，并详细研究了四颗伽利略卫星\(^\text{[177]}\)。

\textbf{伽利略号任务}
\begin{figure}[ht]
\centering
\includegraphics[width=8cm]{./figures/5deadf09cda53153.png}
\caption{1989 年，伽利略号在与火箭对接前的准备阶段} \label{fig_Jupite_12}
\end{figure}
第一艘进入木星轨道的航天器是伽利略号任务，它于 1995 年 12 月 7 日抵达该行星\(^\text{[71]}\)。伽利略号在轨运行超过七年，多次飞掠四颗伽利略卫星及 Amalthea。该航天器还观测到了 1994 年彗星舒梅克–列维 9 号（Shoemaker–Levy 9）与木星相撞的事件。由于高增益天线故障，任务的一些目标未能实现\(^\text{[178]}\)。

一枚 340 公斤的钛制大气探测器于 1995 年 7 月从伽利略号上释放，并于 12 月 7 日进入木星大气\(^\text{[71]}\)。它以约 2,575 公里/小时（1,600 英里/小时）的速度下降，穿越 150 公里（93 英里）的大气层\(^\text{[71]}\)，并在航天器被摧毁前共收集了 57.6 分钟的数据\(^\text{[179]}\)。伽利略轨道器本身于 2003 年 9 月 21 日被故意引导进入木星，经历了更迅速的毁灭。NASA 之所以摧毁伽利略号，是为了避免其可能坠入并污染可能存在生命的木卫二（Europa）\(^\text{[178]}\)。

此次任务的数据表明，氢在木星大气中的含量高达 90\%\(^\text{[71]}\)。记录的温度超过 300 °C（572 °F），在探测器被汽化前测得的风速超过 644 公里/小时（>400 英里/小时）\(^\text{[71]}\)。

\textbf{朱诺号任务}
\begin{figure}[ht]
\centering
\includegraphics[width=8cm]{./figures/3db8b370dca0782e.png}
\caption{2011 年，朱诺号在旋转测试台上准备接受测试} \label{fig_Jupite_13}
\end{figure}
NASA 的朱诺号任务于 2016 年 7 月 4 日抵达木星，目标是在极轨道上对木星进行详细研究。该航天器原本计划在二十个月内绕木星运行三十七圈\(^\text{[78][180][181]}\)。在任务期间，航天器将暴露于来自木星磁层的高强度辐射，这可能导致部分仪器失效\(^\text{[182]}\)。2016 年 8 月 27 日，朱诺号完成了首次飞掠木星，并传回了历史上首批木星北极图像\(^\text{[183]}\)。

朱诺号在 2018 年 7 月预算任务结束前共完成了 12 次轨道运行\(^\text{[184]}\)。当年 6 月，NASA 将任务计划延长至 2021 年 7 月；同年 1 月，任务再次延长至 2025 年 9 月，包含四次卫星飞掠：一次木卫三（Ganymede）、一次木卫二（Europa）以及两次木卫一（Io）\(^\text{[185][186]}\)。当朱诺任务结束时，将进行受控脱轨并在木星大气中解体，以避免撞击并污染木星的卫星\(^\text{[187]}\)。

\textbf{取消的任务与未来计划}

科研界对研究木星较大的冰质卫星非常感兴趣，这些卫星可能拥有地下液态海洋\(^\text{[188]}\)。然而经费困难延缓了任务进展，NASA 的木星冰卫星轨道器（JIMO）计划于 2005 年被取消\(^\text{[189]}\)。随后提出了一项 NASA/ESA 联合任务 EJSM/Laplace，计划于约 2020 年发射。EJSM/Laplace 将由 NASA 主导的木卫二轨道器（Jupiter Europa Orbiter）和 ESA 主导的木卫三轨道器（Jupiter Ganymede Orbiter）组成\(^\text{[190]}\)。但 ESA 于 2011 年 4 月正式终止与 NASA 的合作，原因是 NASA 的预算问题及其对任务时间表的影响。取而代之的是，ESA 决定推进一项完全由欧洲自主执行的任务，以参加其 L1 Cosmic Vision 选拔\(^\text{[191]}\)。这些计划后来实现为欧洲航天局的“木星冰卫星探测器”（JUICE），已于 2023 年 4 月 14 日发射\(^\text{[192]}\)，随后是 NASA 的木卫二探测器 Europa Clipper，于 2024 年 10 月 14 日发射\(^\text{[193]}\)。

其他提出的任务包括中国国家航天局的“天问四号”（Tianwen-4），计划于 2035 年前后向木星系统甚至可能是木卫四（Callisto）发射轨道器\(^\text{[194]}\)，以及 CNSA 的“穿越星际快递”（Interstellar Express）\(^\text{[195]}\)和 NASA 的 Interstellar Probe\(^\text{[196]}\)，两者都将利用木星引力助推抵达日球层边缘。
\subsection{卫星}
木星目前已知拥有 97 颗天然卫星\(^\text{[8]}\)，随着望远镜观测的不断增加，这一数字很可能继续上升\(^\text{[197]}\)。在这些卫星中，仅有 16 颗直径超过 10 公里\(^\text{[198]}\)。其中最大的四颗卫星——伽利略卫星（Galilean moons）——依大小顺序分别为木卫三（Ganymede）、木卫四（Callisto）、木卫一（Io）和木卫二（Europa），在晴朗夜晚使用双筒望远镜即可从地球看到它们\(^\text{[199]}\)。
\subsubsection{伽利略卫星}
由伽利略发现的四颗卫星——Io、Europa、Ganymede 和 Callisto——是太阳系中最大的一批卫星。Io、Europa 与 Ganymede 的轨道构成一种称为拉普拉斯共振（Laplace resonance）的模式：在木卫一完成四次环绕木星的轨道周期时，木卫二正好完成两次，木卫三则正好完成一次。该共振导致三颗大型卫星之间的引力作用在它们每次绕木星的轨道中的同一点对其施加额外拉力，从而将其轨道拉长为椭圆形。另一方面，来自木星的潮汐力则倾向于使它们的轨道趋于圆形\(^\text{[200]}\)。

轨道偏心率使三颗卫星的形状定期发生弯曲：接近木星时，它们被引力拉伸，而远离木星时又恢复为更接近球形。潮汐弯曲产生的摩擦在卫星内部生成热量\(^\text{[201]}\)。这在木卫一上表现得最为显著，因为它受到最强烈的潮汐力\(^\text{[201]}\)；同时，木卫二表面地质的年轻性（这表明其表面曾被近期重新塑造）也在较低程度上体现了这一效应\(^\text{[202]}\)。
\begin{figure}[ht]
\centering
\includegraphics[width=8cm]{./figures/3cf8c93e3bbfd3fa.png}
\caption{} \label{fig_Jupite_14}
\end{figure}
\begin{figure}[ht]
\centering
\includegraphics[width=8cm]{./figures/21f593e7fc447743.png}
\caption{伽利略卫星Io,欧罗巴,木卫三,和木卫四(以增加与木星距离的顺序) 在假色} \label{fig_Jupite_15}
\end{figure}






